\documentclass{beamer}

%-----------------------------------------------------------------
%	Packages
%-----------------------------------------------------------------
\usepackage[utf8]{inputenc}	% für Umlaute ect.
\usepackage{fancyhdr} % für header
\usepackage{lastpage} % für footer
\usepackage{extramarks} % für header und footer
\usepackage{amsthm} % math stuff
\usepackage{amsmath} % math stuff
\usepackage{amssymb} % math stuff
\usepackage{color}
\usepackage{listings} % code listings
\usepackage{graphicx} % für graphics
\usepackage{color}
\usepackage{tikz}
\usepackage[absolute,overlay]{textpos} %to translate graphics through space
\usepackage{soul}
\usepackage{hyperref}


%-----------------------------------------------------------------
%	Title
%-----------------------------------------------------------------
\title{ABS-Normal Form}
\subtitle{Eine Implementierung mit CUDA}
\author[Matthias Mitterreiter]{\includegraphics[height=3cm,width=4cm]{img/badge-nvidia-cuda-cpp}  \\ Matthias Mitterreiter}
\institute{Seminar Parallel Computing - FSU Jena \\
%\vspace{0.5cm}
{\scriptsize Prof. Martin Bücker, Dr. Torsten Bosse, Dipl-Inf. Ralf Seidler}}
\date{\today}

%-----------------------------------------------------------------
%	Settings
%-----------------------------------------------------------------


\mode<presentation>
{
	\usetheme{Warsaw}
	\usecolortheme{crane}
}


%---------------------------------------------------
%	Colors
%---------------------------------------------------

\definecolor{listing-background}{RGB}{77,77,77}
\definecolor{manifest-green}{RGB}{0, 179, 0}
\definecolor{orange}{RGB}{255,127,0}
\definecolor{green}{RGB}{57, 230, 0}
\definecolor{blue}{RGB}{26, 117, 255}
\definecolor{yellow}{RGB}{254, 255, 102}
\definecolor{red}{RGB}{255,77,77}
\definecolor{blue2}{RGB}{51, 204, 255}
\definecolor{lightgray}{rgb}{.9,.9,.9}
\definecolor{darkgray}{rgb}{.4,.4,.4}
\definecolor{purple}{rgb}{0.65, 0.12, 0.82}

%---------------------------------------------------
%	Syntax Highlighting
%---------------------------------------------------

\lstdefinelanguage{manifest}{
  keywords={CACHE, NETWORK, FALLBACK,MANIFEST},
  morecomment=[s]{/*}{*/},
  morecomment=[l]{\#}, 
  morestring=[b]',
  morestring=[b]",
  ndkeywords={class, export, boolean, throw, implements, import, this},
  keywordstyle=\color{manifest-green}\bfseries,
  ndkeywordstyle=\color{black}\bfseries,
  identifierstyle=\color{black},
  commentstyle=\color{purple}\ttfamily,
  stringstyle=\color{red}\ttfamily,
  sensitive=true
}

\lstdefinelanguage{JavaScript}{
  keywords={CACHE, NETWORK, FALLBACK,MANIFEST, break, case, catch, continue, debugger, default, delete, do, else, false, finally, for, function, if, in, instanceof, new, null, return, switch, this, throw, true, try, typeof, var, void, while, with},
  morecomment=[l]{//},
  morecomment=[s]{/*}{*/},
  morestring=[b]',
  morestring=[b]",
  ndkeywords={class, export, boolean, throw, implements, import, this},
  keywordstyle=\color{blue}\bfseries,
  ndkeywordstyle=\color{darkgray}\bfseries,
  identifierstyle=\color{black},
  commentstyle=\color{purple}\ttfamily,
  stringstyle=\color{red}\ttfamily,
  sensitive=true
}

\lstset{
   language=JavaScript,
   backgroundcolor=\color{lightgray},
   extendedchars=true,
   basicstyle=\footnotesize\ttfamily,
   showstringspaces=false,
   showspaces=false,
   numbers=left,
   numberstyle=\footnotesize,
   numbersep=7pt,
   tabsize=2,
   breaklines=true,
   showtabs=false,
   captionpos=b,
   literate={\ \ }{{\ }}1
}

% An environment for stpes, cases ect. 
% From: http://tex.stackexchange.com/questions/32798/a-step-by-step-environment
\newenvironment{steps}[1]{\begin{enumerate}[label=#1 \arabic*]}{\end{enumerate}}
\makeatletter%
\def\step{%
	\@ifnextchar[ \@step{\@noitemargtrue\@step[\@itemlabel]}}
\def\@step[#1]{\item[#1]\mbox{}\\\hspace*{\dimexpr-\labelwidth-\labelsep}}
\makeatother

%-----------------------------------------------------------------------------
%	Theoreme
%-----------------------------------------------------------------------------
\newtheorem{mydef}{Definition}
\newtheorem*{mydef*}{Definition}
\newtheorem{mybei}{Beispiel}
\newtheorem*{mybei*}{Beispiel}\usepackage{tikz}
\newtheorem{mysatz}{Satz}
\newtheorem*{mysatz*}{Satz}
\newtheorem{mybew}{Beweis}
\newtheorem*{mybew*}{Beweis}
\newtheorem{myfolg}{Folgerung}
\newtheorem*{myfolg*}{Folgerung}
\newtheorem{mybemerk}{Bemerkung}
\newtheorem*{mybemerk*}{Bemerkung}


%-------------------------------------------------
%	Beamer
%-------------------------------------------------
%\definecolor{craneorange}{rgb}{0.68,1,1}
%\definecolor{craneblue}{gray}{0.85}
\begin{document}

\maketitle

\addtobeamertemplate{frametitle}{}{%
\begin{tikzpicture}[remember picture,overlay]
\node[anchor=north east,yshift=1pt] at (current page.north east) {\includegraphics[height=0.8cm]{img/badge-nvidia-cuda-cpp}};
\end{tikzpicture}}
	\begin{frame}
\frametitle{Inhalt}
\framesubtitle{...}
	\begin{enumerate}
		\item ABS - Normal Form
		\item Aufgabengeschreibung
		\item Implementierung
		\begin{enumerate}
			\item Verwendete Bibliotheken
			\item Evaluate
			\item Gradient
			\item Solve
		\end{enumerate}
		\item Anwendung
		\item Numerische Tests und Vergleiche
		\item Cool Snippets
		\item Imprevements
		\item Problems Lesson Learned
	\end{enumerate}
\end{frame}
%------------------------------------------------------------------------
\begin{frame}
	\frametitle{ABS-NF}
	\framesubtitle{Einführung}
	\setbeamertemplate{enumerate items}[default]
	\begin{mydef*}
		Smooth funcion \\
		A smooth function $f(x)$ is continuous and has a continuous derivative
	\end{mydef*}
	$f(x) = x^2$ $f'(x) = 2x$

\end{frame}
%------------------------------------------------------------------------
\begin{frame}
	\frametitle{ABS-NF}
	\framesubtitle{Einführung}
	\setbeamertemplate{enumerate items}[default]
	\begin{mydef*}
		Picewise smooth function \\
		A picewise smooth function $f(x)$ is made up of finitly many smooth functions, separated by
		jump discontiuities.
	\end{mydef*}
	\begin{flalign*}
		f(x) = \begin{cases}
			1 & x < -1 \\
			x^2 + 1 & -1 <= x < 2 \\
			x & 2 <= x
		\end{cases}
	\end{flalign*}
	Wir betrachten ausschließlich continuous picewise smooth functions.
\end{frame}
%------------------------------------------------------------------------
\begin{frame}
\frametitle{ABS-NF}
\framesubtitle{Einführung}
	Picewise smooth  funcitons können durch picewise linear functions (PL) approximiert werden.
	\begin{center}
		BILD
	\end{center}
	Betrachten ausschließlich continuous PL
\end{frame}
%------------------------------------------------------------------------
\begin{frame}
	\frametitle{ABS-NF}
	\framesubtitle{Einführung}
	\setbeamertemplate{enumerate items}[default]
	ABS-Normal Form:
	\begin{center}
		Repräsentierung für Picewise linear functions (PL) 
	\end{center}
	\begin{flalign*}
	\begin{pmatrix}
	\Delta z \\
	\Delta y
	\end{pmatrix}
	= 
	\begin{pmatrix}
	a \\
	b
	\end{pmatrix}
	+
	\begin{pmatrix}
	Z & L \\
	J & Y 
	\end{pmatrix}
	\times
	\begin{pmatrix}
	\Delta x \\
	|\Delta z |
	\end{pmatrix}
	\end{flalign*}
\end{frame}
%------------------------------------------------------------------------
\begin{frame}
	\frametitle{ABS-NF}
	\framesubtitle{Einführung}
	\setbeamertemplate{enumerate items}[default]
	Vorgehen:
	\begin{itemize}
		\item Any picewise linear scalar function has a so called min-max repräsentation
		\item All min max expressions can be expressed in terms of abs functions
		\item Picewise linearization wird erreicht durch algorithmic differentiation
	\end{itemize}
\end{frame}
%------------------------------------------------------------------------
\begin{frame}
	\frametitle{ABS-NF}
	\framesubtitle{Beispiel}
	\setbeamertemplate{enumerate items}[default]
	\begin{flalign*}
		F(x_1,x_2) &= (x_2^2 - x_1^+)^+ \\
			 (i)^+ &= \max(0, i)
	\end{flalign*}
	\begin{center}
		BILD
	\end{center}
	
\end{frame}
%------------------------------------------------------------------------
\begin{frame}
	\frametitle{ABS-NF}
	\framesubtitle{Einführung}
	\setbeamertemplate{enumerate items}[default]
	\begin{center}
		ADOL C, Beispiel
	\end{center}
\end{frame}
%------------------------------------------------------------------------
\begin{frame}
	\frametitle{ABS-NF}
	\framesubtitle{Aufgabenbeschreibung}
	\setbeamertemplate{enumerate items}[default]
	\begin{itemize}
		\item Evaluate abs normal form
		\item Calculate Gradient
		\item Solve abs-normal form system
	\end{itemize}
\end{frame}
%------------------------------------------------------------------------
	\begin{frame}
\frametitle{Problems}
\framesubtitle{...}
	\begin{enumerate}
		\item ABS - Normal Form
		\item C/C++
		\item CUDA
	\end{enumerate}
\end{frame}
%------------------------------------------------------------------------
\end{document}